\scanpage{42}
fibré correspondant l'ens de ses sections, qui est un torseur sur le groupe \inserted{$Z$} des sections du fibré vectoriel correspondant, lequel est un $\mathbb Z/n\mathbb Z$-module libre de rang 1, i.e. un groupe cyclique d'ordre $n$.

Si $S$ est un quasi-pol. combinatoire, $Z$ son groupe des rotations, le groupe de ses automorphismes \inserted{(noté $\mathrm{Aff}(S)$)} est canoniquement muni d'une structure d'extension
\diagram{\begin{tikzcd}[ampersand replacement=\&,column sep=small,row sep=small]
0\arrow[r] \& Z\arrow[r] \& \mathrm{Aff}(S)\arrow[r] \& (\mathbb Z/n\mathbb Z)^{*}\arrow[r]\arrow[d,phantom,"\rotatebox{90}{$\simeq$}" description] \& 0\\
\& \& \& \operatorname{Aut}(Z) \&
\end{tikzcd}}
et le choix d'un $s\in S$ détermine un splitting de cette extension, permettant d'identifier $\mathrm{Aff}(S)$ au produit $\frac12$ direct $(\mathbb Z/n\mathbb Z)^*\mathbin{.}Z$. L'image inverse du ss-grpe $\pm1$ de $(\mathbb Z/n\mathbb Z)^*$ dans $\mathrm{Aff}(S)$ est un ss-grpe $\mathrm{Aff}^{\pm}(S)$ de $\mathrm{Aff}(S)$ d'ordre $2n$. C'est aussi le
% margin note written sideways
\margin{ss-groupe $\operatorname{Aut}(S,\omega)$ de $\unsure{\mathfrak S_S|\mathcal S_S|G_S}$, où $\omega$ est n'importe laquelle des str. pol. sur $S$ compatibles avec sa str. quasi-pol.}

Soit $E$ un plan affine sur corps $k$.

Un \uline{système régulier de sommets} dans $E$ est \deleted{\unsure{la co}} un couple $(S,Z)$, où $S\subset E$ est une partie finie,
% archive sheet number: 41
